\documentclass[preview]{standalone}

\usepackage{amsmath}
\usepackage{amssymb}
\usepackage{stellar}
\usepackage{definitions}
\usepackage{bettelini}

\begin{document}

\id{surface-integrals}
\genpage

\section{Surface Integrals}

\begin{snippetdefinition}{scalar-surface-integral-definition}{Scalar surface integral}
    Let
    \(\sigma\colon\overline A\subseteq\realnumbers^2
    \fromto\realnumbers^3\) be a regular simple parametrization of
    \(\Sigma=\sigma(\overline A)\), and let
    \(f\colon\Sigma\fromto\realnumbers\) be continuous. Define
    \[
        \iint_\Sigma f\,dS
        =
        \iint_A
        f(\sigma(u,v))
        \|\sigma_u(u,v)\times\sigma_v(u,v)\|\,du\,dv.
    \]
    In particular,
    \[
        \operatorname{Area}(\Sigma)
        =\iint_A\|\sigma_u\times\sigma_v\|\,du\,dv.
    \]
\end{snippetdefinition}

\begin{snippet}{surface-area-element-interpretation}
    The vector \(\sigma_u\times\sigma_v\) is perpendicular to the tangent
    plane. Its norm is the infinitesimal area-dilation factor, so
    \[
        dS=\|\sigma_u\times\sigma_v\|\,du\,dv.
    \]
    Measure-theoretically, this is the area measure induced on the image by
    the parametrization; the factor records its local deformation.
\end{snippet}

\begin{snippetdefinition}{parametrized-surface-unit-normal-definition}{Unit normal associated with a parametrization}
    Let
    \(\sigma\colon\overline A\subseteq\realnumbers^2
    \fromto\realnumbers^3\) be regular. The associated unit normal is
    \[
        \nu(u,v)=
        \frac{\sigma_u(u,v)\times\sigma_v(u,v)}
        {\|\sigma_u(u,v)\times\sigma_v(u,v)\|}.
    \]
    Reversing the order of the parameters replaces \(\nu\) by \(-\nu\).
\end{snippetdefinition}

\begin{snippetexample}{graph-surface-area-element-example}{Area element of a graph}
    For \(\sigma(x,y)=(x,y,f(x,y))\),
    \[
        \sigma_x=(1,0,f_x),
        \qquad
        \sigma_y=(0,1,f_y),
    \]
    and hence
    \[
        \sigma_x\times\sigma_y=(-f_x,-f_y,1),
        \qquad
        dS=\sqrt{1+\|\nabla f\|^2}\,dx\,dy.
    \]
\end{snippetexample}

\section{Classical Surfaces}

\begin{snippetexample}{sphere-surface-parametrization-example}{Sphere}
    The upper hemisphere of the unit sphere can first be represented as the
    graph
    \[
        \sigma(x,y)=\left(x,y,\sqrt{1-x^2-y^2}\right),
        \qquad x^2+y^2<1.
    \]
    Replacing the square root by its negative gives the lower hemisphere.
    A spherical-coordinate parametrization of the sphere of radius \(R\) is
    \[
        \sigma(\varphi,\theta)=
        (R\sin\varphi\cos\theta,
         R\sin\varphi\sin\theta,
         R\cos\varphi),
    \]
    where \(\varphi\in[0,\pi]\) and
    \(\theta\in[0,2\pi]\). Its coordinate vectors are
    \[
        \sigma_\varphi=
        (R\cos\varphi\cos\theta,
         R\cos\varphi\sin\theta,
         -R\sin\varphi)
    \]
    and
    \[
        \sigma_\theta=
        (-R\sin\varphi\sin\theta,
         R\sin\varphi\cos\theta,
         0).
    \]
    Thus
    \[
        \sigma_\varphi\times\sigma_\theta
        =R^2\sin\varphi
        (\sin\varphi\cos\theta,
         \sin\varphi\sin\theta,
         \cos\varphi),
    \]
    and therefore
    \[
        \|\sigma_\varphi\times\sigma_\theta\|
        =R^2\sin\varphi.
    \]
    The parametrization is regular for \(0<\varphi<\pi\), but degenerates
    at the poles; the parameter values \(\theta=0\) and \(\theta=2\pi\)
    are also identified along one meridian. Its normal is outward: at
    \((\varphi,\theta)=(\pi/2,0)\), both the position and normal directions
    point along the positive \(x\)-axis.
\end{snippetexample}

\begin{snippetcorollary}{sphere-surface-area-corollary}{Area of a sphere}
    The sphere of radius \(R\) has area
    \[
        \operatorname{Area}(S_R^2)
        =\integral[0][2\pi]
        [\integral[0][\pi][R^2\sin\varphi][\varphi]][\theta]
        =4\pi R^2.
    \]
\end{snippetcorollary}

\begin{snippetexample}{cylinder-surface-area-element-example}{Cylinder}
    The cylinder \(x^2+y^2=R^2\) is parametrized by
    \[
        \sigma(\theta,z)=(R\cos\theta,R\sin\theta,z).
    \]
    Since
    \[
        \sigma_\theta\times\sigma_z
        =(R\cos\theta,R\sin\theta,0),
    \]
    the normal is outward and \(dS=R\,d\theta\,dz\). For \(R=1\), the
    parametrization preserves area locally. The values at \(\theta=0\) and
    \(\theta=2\pi\) are identified along the boundary of the parameter
    domain.
\end{snippetexample}

\begin{snippetexample}{torus-surface-area-example}{Torus}
    Let \(R>r>0\). Rotating the circle
    \(\gamma(\theta)=(R+r\cos\theta,r\sin\theta)\) about the \(z\)-axis
    gives
    \[
        \sigma(\theta,\varphi)=
        \bigl((R+r\cos\theta)\cos\varphi,
              (R+r\cos\theta)\sin\varphi,
              r\sin\theta\bigr).
    \]
    Its area factor is
    \[
        \|\sigma_\theta\times\sigma_\varphi\|
        =r(R+r\cos\theta)>0.
    \]
    Therefore
    \begin{align*}
        \operatorname{Area}(T)
        &=\iint_{[0,2\pi]^2}r(R+r\cos\theta)
            \,d\theta\,d\varphi\\
        &=4\pi^2rR.
    \end{align*}
    This is the product of the length \(2\pi r\) of the generating circle
    and the distance \(2\pi R\) traveled by its centroid, in agreement with
    Pappus's theorem.
\end{snippetexample}

\end{document}
